\documentclass[11pt,a4paper]{article}
\usepackage[margin=0.85in,top=1in,bottom=1in]{geometry}
\usepackage{amsmath,amssymb,amsfonts}
\usepackage{graphicx}
\usepackage{float}
\usepackage{enumitem}
\usepackage{microtype}
\usepackage{caption}
\usepackage[hidelinks]{hyperref}
\usepackage[most,breakable]{tcolorbox}
\tcbuselibrary{breakable,skins}
\usepackage{fontspec}
\setmainfont{DejaVu Serif}
\setsansfont{DejaVu Sans}
\setmonofont{DejaVu Sans Mono}
\captionsetup{font=small,labelfont=bf,skip=4pt}
\setlist{leftmargin=1.5em,topsep=2pt}

%% Breakable card — prevents content cutoff across pages
\newtcolorbox{solcard}[3]{
  breakable, enhanced,
  colback=white, colframe=gray!50, arc=2pt,
  title={\textbf{#1}\hfill\textsf{\small #3\,pts}},
  fonttitle=\normalsize\sffamily\bfseries,
  before upper={\small\textit{#2}\par\smallskip\hrule height 0.4pt\smallskip},
  top=4pt, bottom=6pt, left=7pt, right=7pt,
  parbox=false,
}

\title{\textbf{Mixture viscosity measurement --- Solution}}
\author{\normalsize W26 (2026) — English translation}
\date{}
\begin{document}\maketitle\vspace{-0.5em}
\begin{solcard}{A1}{Express the molar concentration of water $n_w$ in solution in terms of the mass concentration of glycerol $w_g$, $\mu_w$ and $\mu_g$.}{0.50}

The amount of a substance in a mixture is defined as: \[\nu=\dfrac{m}{\mu},\]where $m$ is the mass of the substance, and $\mu$ is its molar mass. Let's write the formula for the molar concentration of water: \[n_w=1-n_g=1-\dfrac{\nu_g}{\nu_g+\nu_w},\]from here:

\par\smallskip\noindent\textbf{Answer:} \textbf{Answer:} \[n_w=1-\frac{w_g\mu_w}{\mu_g-w_g(\mu_g-\mu_w)}\]
\end{solcard}

\begin{solcard}{A2}{Using the equation $(2)$, convert the values ​​of $T_g$ to the values ​​of $w_g$.}{2.00}

Using the table method or the method of simple iterations with $a=0.9$, we get:

\par\smallskip\noindent\textbf{Answer:} \textbf{Answer:} $T_g,~\mathrm{K}$ $w_g$ $\eta,~\mathrm{Pa}\cdot\mathrm{s}$ 195.2 1.000 0.906 194.8 0.996 0.839 194.6 0.994 0.767 193.9 0.987 0.655 193.1 0.979 0.558 191.8 0.965 0.479 190.8 0.954 0.373 188.6 0.930 0.202 183.4 0.869 0.0674 173.0 0.728 0.0278 170.0 0.682 0.0134 154.6 0.384 0.0026 149.3 0.252 0.0017 143.8 0.100 0.0012 140.3 0 0.0009
\end{solcard}

\begin{solcard}{A3}{Propose a linearization that describes a system with a low molar concentration of water, which can be used to determine the parameters $\xi$ and $k$.}{0.30}
\par\smallskip\noindent\textbf{Answer:} \textbf{Answer:} The dependence $\operatorname{ln}\dfrac{\eta}{1~\mathrm{Pa}\cdot \mathrm{s}}(n_w)$ is linear.
\end{solcard}

\begin{solcard}{A4}{Plot the linearized relationship over the range $n_w \in [0, 0.2]$. Determine the parameters $\xi$ and $k$ using the resulting graph.}{1.20}

Let's recalculate the points into the dependence $\operatorname{ln}\dfrac{\eta}{1~\mathrm{Pa}\cdot \mathrm{s}}(n_w)$.


$T_g,~\mathrm{K}$ $w_g$ $\eta,~\mathrm{Pa}\cdot\mathrm{s}$ $n_w$ $\operatorname{ln} \eta$ 195.2 1.000 0.906 0 -0.099 194.8 0.996 0.839 0.020 -0.176 194.6 0.994 0.767 0.030 -0.265 193.9 0.987 0.655 0.063 -0.423 193.1 0.979 0.558 0.099 -0.583 191.8 0.965 0.479 0.156 -0.736 190.8 0.954 0.373 0.198 -0.986 188.6 0.930 0.202 0.278 -1.600 183.4 0.869 0.0674 173.0 0.728 0.0278 170.0 0.682 0.0134 154.6 0.384 0.0026 149.3 0.252 0.0017 143.8 0.100 0.0012 140.3 0 0.0009


$\operatorname{ln}\dfrac{\eta}{1~\mathrm{Pa}\cdot \mathrm{s}}(n_w)$


Approximating the dependence of the straight line, we obtain: \[k=4.31,\] \[\operatorname{ln}\dfrac{\xi}{1~\mathrm{Pa}\cdot \mathrm{s}}=-0.12,\]from which it follows:

\par\smallskip\noindent\textbf{Answer:} \textbf{Answer:} \[k=4.31\]\[\xi=0.887~\mathrm{Pa}\cdot \mathrm{s}\]
\begin{figure}[H]
  \centering
  \includegraphics[width=0.72\textwidth]{PE_S_EN_assets/fig1.jpg}
\end{figure}
\end{solcard}

\begin{solcard}{A5}{Knowing the parameters $\xi$ and $k$, propose a linearization that allows you to determine the values ​​of the parameters $a$ and $b$.}{0.30}
\par\smallskip\noindent\textbf{Answer:} \textbf{Answer:} The dependence $\operatorname{ln}|\dfrac{\xi \operatorname{exp}(-kn_w)}{\eta}-1|(n_w)$ is linear.
\end{solcard}

\begin{solcard}{A6}{Plot the linearized relationship. From the graph, determine $a$ and $b$. Note: due to the inaccuracy of experimental data, some values ​​cannot be recalculated. There are no more than 5 such values.}{1.20}

Note: due to the inaccuracy of experimental data, some values ​​cannot be recalculated. There are no more than 5 such values.


Let's continue recalculating points:


$T_g,~\mathrm{K}$ $w_g$ $\eta,~\mathrm{Pa}\cdot\mathrm{s}$ $n_w$ $\operatorname{ln}|\dfrac{\xi \operatorname{exp}(-kn_w)}{\eta}-1|$ 195.2 1.000 0.906 0 -3.92 194.8 0.996 0.839 0.020 -3.54 194.6 0.994 0.767 0.030 -4.12 193.9 0.987 0.655 0.063 -3.44 193.1 0.979 0.558 0.099 -3.28 191.8 0.965 0.479 0.156 -2.91 190.8 0.954 0.373 0.198 -4.35 188.6 0.930 0.202 0.278 -1.12 183.4 0.869 0.0674 0.435 0.02 173.0 0.728 0.0278 0.656 -0.12 170.0 0.682 0.0134 0.704 0.78 154.6 0.384 0.0026 0.891 1.85 149.3 0.252 0.0017 0.938 2.10 143.8 0.100 0.0012 0.979 2.29 140.3 0 0.0009 1 2.50


Let's plot the resulting points on the graph:


$\operatorname{ln}|\dfrac{\xi \operatorname{exp}(-kn_w)}{\eta}-1|(n_w)$


In the range $n_w \in [0, 0.2]$, the expression under the logarithm has a different sign and is near zero, so the function in this region behaves non-monotonic. To draw a straight line, we take into account only points in the range $n_w \in [0.7, 1].$


$\operatorname{ln}|\dfrac{\xi \operatorname{exp}(-kn_w)}{\eta}-1|(n_w)$ in the range $n_w \in [0.7, 1]$ (not required from participants)


For $n_w>0.7$, $\dfrac{\xi \operatorname{exp}(-kn_w)}{\eta}>1$ holds, so $a>0$. Having drawn the approximating straight line, we find: \[b=5.66,\]\[\operatorname {ln} a=-3.21,\]from here:

\par\smallskip\noindent\textbf{Answer:} \textbf{Answer:} \[a=0.040\]\[b=5.66\]
\begin{figure}[H]
  \centering
  \includegraphics[width=0.72\textwidth]{PE_S_EN_assets/fig2.jpg}
\end{figure}
\begin{figure}[H]
  \centering
  \includegraphics[width=0.72\textwidth]{PE_S_EN_assets/fig3.jpg}
\end{figure}
\end{solcard}

\begin{solcard}{A7}{Draw a graph of the theoretical and experimental dependences $\eta(n_w)$ on one piece of graph paper.}{1.00}

Let's calculate the theoretical values ​​of $\eta(n_w)$:


$T_g,~\mathrm{K}$ $w_g$ $\eta,~\mathrm{Pa}\cdot\mathrm{s}$ $n_w$ $\operatorname{ln} \eta$ $\operatorname{ln}|\dfrac{\xi \operatorname{exp}(-kn_w)}{\eta}-1|$ $\eta_{теор},~\mathrm{Pa}\cdot\mathrm{s}$ 195.2 1.000 0.906 0 -0.099 -3.92 0.853 194.8 0.996 0.839 0.020 -0.176 -3.54 0.779 194.6 0.994 0.767 0.030 -0.265 -4.12 0.744 193.9 0.987 0.655 0.063 -0.423 -3.44 0.640 193.1 0.979 0.558 0.099 -0.583 -3.28 0.541 191.8 0.965 0.479 0.156 -0.736 -2.91 0.413 190.8 0.954 0.373 0.198 -0.986 -4.35 0.337 188.6 0.930 0.202 0.278 -1.600 -1.12 0.224 183.4 0.869 0.0674 0.435 0.02 0.0926 173.0 0.728 0.0278 0.656 -0.12 0.0199 170.0 0.682 0.0134 0.704 0.78 0.0135 154.6 0.384 0.0026 0.891 1.85 0.0026 149.3 0.252 0.0017 0.938 2.10 0.0017 143.8 0.100 0.0012 0.979 2.29 0.0012 140.3 0 0.0009 1 2.50 0.0010


Let's build a graph of the obtained dependencies:


$\eta(n_w)$


Having drawn smoothing curves, we conclude: the equation is applicable over the entire range ($n_w \in [0, 1]$).

\begin{figure}[H]
  \centering
  \includegraphics[width=0.72\textwidth]{PE_S_EN_assets/fig4.jpg}
\end{figure}
\end{solcard}

\begin{solcard}{B1}{Using the equation $(2)$, determine the glass transition temperature of the solution $T_g$.}{0.50}

Let's express $w_g$: \[w_g=\dfrac{m_g}{m_g+m_w}=\dfrac{\rho_gc_g}{(\rho_g-\rho_w)c_g+\rho_w},\] then:

\par\smallskip\noindent\textbf{Answer:} \textbf{Answer:} $T_g=\left(140.3+35.272\dfrac{\rho_gc_g}{(\rho_g-\rho_w)c_g+\rho_w}-3.879\left(\dfrac{\rho_gc_g}{(\rho_g-\rho_w)c_g+\rho_w}\right)^2+23.467\left(\dfrac{\rho_gc_g}{(\rho_g-\rho_w)c_g+\rho_w}\right)^3\right)~\mathrm{K}\approx 162.8~\mathrm{K}$
\end{solcard}

\begin{solcard}{B2}{Propose a linearization to determine the values ​​of the parameters $D$ and $T_0$. After plotting the linearized dependence, determine the values ​​of these parameters.}{2.00}

The dependence $\dfrac{1}{\operatorname {ln}\dfrac{\eta}{\eta_0}}(T)$ is linear with slope $k=\dfrac{1}{DT_0}$ and intercept term $b=-\dfrac{1}{D}$. Let’s recalculate the point using this formula and plot the linearized dependence:


$T,~^\circ C$ $\eta,~\mathrm{Pa}\cdot\mathrm{s}$ $T,~\mathrm{K}$ $\dfrac{1}{\operatorname {ln}\dfrac{\eta}{\eta_0}}$ 1 0.0201 274 0.556 4 0.0173 277 0.607 11 0.0128 284 0.743 15 0.0104 288 0.878 35 0.0046 308 3.10 44 0.0034 317 48.1 50 0.0031 323 -14.0 56 0.0026 329 -4.04 60 0.0022 333 -2.41 67 0.0020 340 -1.96 71 0.0017 344 -1.487 76 0.00145 349 -1.20 85 0.0013 358 -1.06 90 0.0011 363 -0.903 96 0.0010 369 -0.831 100 0.0009 373 -0.764


Let's plot the resulting relationship:


$\dfrac{1}{\operatorname {ln}\dfrac{\eta}{\eta_0}}(T)$


From the proposed equation it follows that it can only be fulfilled at sufficiently low temperatures. It is clear from the graph that the equation is applicable for the first 4 points of the graph (at $T<288~\mathrm{K}$). Let's demonstrate this clearly:


$\dfrac{1}{\operatorname {ln}\dfrac{\eta}{\eta_0}}(T)$ at $T<288~K$ (not required from participants)


Having determined the parameters of the approximating straight line, we obtain a system of equations: \[\begin{cases}-\dfrac{1}{D}=-5.60\\ \dfrac{1}{DT_0}=0.022 К^{-1},\end{cases}\] which follows:

\par\smallskip\noindent\textbf{Answer:} \textbf{Answer:} \[T_0=255~\mathrm{K}\]\[D=0.18\]
\begin{figure}[H]
  \centering
  \includegraphics[width=0.72\textwidth]{PE_S_EN_assets/fig5.jpg}
\end{figure}
\begin{figure}[H]
  \centering
  \includegraphics[width=0.72\textwidth]{PE_S_EN_assets/fig6.jpg}
\end{figure}
\end{solcard}

\begin{solcard}{С1}{Determine the viscosity of the mixture $\eta_n$ at the volume concentration of glycerol $c_g=0.5$ and temperature $T=25~^\circ \mathrm C$ using the dependence on the molar concentration of water.}{0.40}

The equation obtained for $\eta_n(n_w)$ is applicable over the entire range. Therefore, it is enough to convert $c_g$ into $n_w$: \[n_w=1-\frac{w_g\mu_w}{\mu_g-w_g(\mu_g-\mu_w)}=1-\frac{\dfrac{\rho_gc_g}{(\rho_g-\rho_w)c_g+\rho_w}\mu_w}{\mu_g-\dfrac{\rho_gc_g}{(\rho_g-\rho_w)c_g+\rho_w}(\mu_g-\mu_w)}\approx 0.80\] and substitute it into the resulting equation. Thus we get:

\par\smallskip\noindent\textbf{Answer:} \textbf{Answer:} $\eta_n=6.0~мПа\cdot \mathrm{s}$
\end{solcard}

\begin{solcard}{С2}{Determine the viscosity of the mixture $\eta_t$ under the same conditions using the temperature dependence.}{0.40}

The equation obtained for $\eta_t(T)$ is indispensable for $T=25^\circ C$. You can determine the viscosity at a given temperature, for example, graphically by plotting the curve $\eta_t(T)$, but the answer with good accuracy is given by a linear approximation: \[\eta_{25}\approx\dfrac{\eta_{15}+\eta_{35}}{2}=7.5~мПа\cdot \mathrm{s}.\]

\par\smallskip\noindent\textbf{Answer:} \textbf{Answer:} \[\eta_t=7.5~мПа\cdot \mathrm{s}\]
\end{solcard}

\begin{solcard}{С3}{Compare your results. If there are discrepancies, suggest a possible reason for them.}{0.20}
\par\smallskip\noindent\textbf{Answer:} \textbf{Answer:} The meanings of $\eta_t$ and $\eta_n$ are divergent. One of the main possible reasons for discrepancies in results is poorly measured temperature (for example, if the room temperature changed during the experiment), because the dependence of the viscosity of the mixture on temperature is very strong.
\end{solcard}

\end{document}